\documentclass[a4paper,11pt,fleqn]{scrartcl}%fleqn pour aligner les equations à gauche
\usepackage{../../Styles/Cours}


\newcommand{\chnum}{Chapitre 15}
\newcommand{\chtitle}{Les ondes mécaniques}
\newcommand{\dtype}{Cours}

\begin{document}
	\maketitle
 On propose dans ce TP de mesurer la vitesse du son dans l'air.
\begin{materiel}
    \begin{itemize}
        \item un émetteur et deux récepteurs d'ondes ultrasonores
        \item une réglette
        \item une platine d'acquisition
        \item un générateur
        \item des cables de branchements
        \item un ordinateur équipé de LatisPro
    \end{itemize}
\end{materiel}

\section{Mesure de la célérité d'une onde sonore}
\begin{protocole}
    \textbf{Branchements : }
\begin{itemize}
    \item L'émetteur à ultrasons $E$ est alimenté avec une tension continue de \SI{15}{V}.
	\item Le récepteur à ultrasons transforme le signal perçu en signal électrique de même fréquence. On visualise les variations de ce signal électrique sur l'ordinateur.
	\item Sélectionner le mode continu sur l'émetteur $E$
	\item Pour pouvoir enregistrer les signaux reçus, relier le récepteur 1 en $EA0$, le récepteur 2 en $EA1$ (ne pas oublier de relier les masses en noir à celle de la carte d'acquisition)
	\item Aligner l'émetteur et les deux récepteurs le long d'un rail gradué.
\end{itemize}
    \textbf{Réglages de LatisPro : }
    \begin{itemize}
        \item Ouvrir le logiciel LatisPro
		\item Cliquer sur les voies $EA0$ et $EA1$ pour les activer
		\item Changer le style de courbe. Pour cela, faire un clic-droit sur la légende $EA1$ et cliquer sur "Propriétés" puis "Style" et choisir "Trait". Valider "Ok". N'oublier pas de faire la même chose avec $EA0$.
    \end{itemize}
    \textbf{Mesures : }
    \begin{itemize}
        \item Décaler  le récepteur A d'une distance $d$ suffisamment grande pour pouvoir mesurer avec précision le retard ultrasonore $\tau$ entre les deux récepteurs. Noter les valeurs de $d$ et $\tau$ dans le tableau ci-après.
        \item Répéter l'opération pour avoir 10 mesures différentes.
    \end{itemize}
\end{protocole}

\begin{enumerate}
    \item Représenter le dispositif expérimental et dessiner l'écran  affichant les courbes. Légender puis indiquer avec une double flèche comme se pesurent $\tau$ et $d$ sur ces schémas.
    \item à partir des valeurs obtenues expérimentalement (couples de valeurs $d$ et $\tau$), calculer la valeur de la célérité de l'onde ultrasonore. Compléter le tableau suivant.
    \item À partir des valeurs obtenues expérimentalement (couples de valeurs $d$ et $\Delta t$) : \begin{enumerate}
        \item Tracer le graphe montrant l'évolution de la grandeur $d$ en fonction de $\tau$.
        \item Choisir un modèle mathématique pertinent pour modéliser ces données.
        \item Faire apparaitre l'équation de la courbe de tendance ainsi que son coefficient de corrélation ($R^2$). pour que le modèle soit valide il faut que $R^2>=0,98$.
    \end{enumerate}
    \item Déterminer la valeur de la célérité des ultrasons dans l'air à partir du modèle obtenur précédemment. Pour cela utiliser l'équation obtenur et expliquer votre démarche.
    \item Calcuelr l'écart relatif $e$ sur la célérité, sachant que la célérité des ultrasons dans l'aur à une température de \SI{20}{\degree} est : $c_{ref} = \SI{340}{\meter\per\second}$. \\
    L'écart relatif $e$ est donné par la relation : $ee =\dfrac{|c_{exp} - c_{ref}|}{c_{ref}}$ où $c_{ref}$ est la célérité déterminée par l'expérience tandis que $c_{ref}$ est la valeur de référence.
\end{enumerate}

\section{Détermination des caractéristiques d'une onde périodique}
\subsection{Période et fréquence}
\begin{protocole}
    \begin{itemize}
        \item Utiliser l'émetteur en mode continu et un seul récepteur proche de l'émetteur.
        \item Déterminer la période $T$ par lecture graphique.
    \end{itemize}
\end{protocole}
\begin{enumerate}[resume]
    \item Représenter le signal obtenu en faisant apparaitre une période (en rouge).
    \item Noter la valeu de la période $T$ obtenue par lecture graphique et la convertir en seconde.
    \item Calculer la fréquence du signal.
\end{enumerate}

\subsection{Longueur d'onde}
\begin{protocole}
    \begin{itemize}
        \item Utiliser deux récepteurs en les plaçant à la même distance de l'émetteur, en mode continu.
        \item Faire les réglages sur longitudinalesisPro de manière à ce que les deux signaux apparaissent, soient discernables et parfaitement superposés.
        \item Les deux signaux sont en phase, c'est à dire que leurs maximums et leurs minimums coïncident.
        \item Déplacer lentement l'un des deux récepteurs de manière à ce que les deux signaux se décalent jusqu'à être de nouveau en phase.
        \item Mesurer la distance entre les deux récepteurs, elle correspond à la longueur d'onde $\lambda$.
    \end{itemize}
\end{protocole}

\begin{enumerate}[resume]
    \item Déterminer la valeur de la longueur d'onde $\lambda$ de la manière la plus précise possible. noter cette valeur avec son unité en déteillant la méthode.
    \item Vérifier la cohérence des résultats obtenus à partir des mesures, en calculant la célérité $c$ de l'onde à partir de la fréquence $f$ et la longueur d'onde $\lambda$.
    \item Commenter le résultat obtenu.
\end{enumerate}
\end{document}